%% chapters/01-getting-started.tex
\chapter{Getting Started}\label{chap:getting-started}

This template produces a landscape A4 book with a professional look and a
small footprint. This chapter explains how the project is organised, how to
compile it and how to add your own chapters.

\section{What You Get}

\begin{itemize}
  \item A landscape A4 page set in two columns, or in one column with a wide
        margin for notes (see \cref{sec:layout}).
  \item Front matter, numbered chapters, appendices and a bibliography.
  \item Callout boxes, theorem-like environments and code listings.
  \item Clickable cross-references and a linked table of contents.
  \item Only widely available packages, so the project compiles quickly and
        works on any \TeX{} distribution.
\end{itemize}

\section{Project Structure}\label{sec:structure}

The project is split into small files, each with a single purpose
(\cref{lst:structure}). Keeping the content of each chapter in its own file
makes the book easier to navigate, review and version-control.

\begin{lstlisting}[language={}, caption={Layout of the project.}, label={lst:structure}, float=tb]
.
|-- main.tex            metadata and structure
|-- references.bib      bibliography database
|-- style/
|   `-- booktemplate.sty  layout, fonts, packages
|-- frontmatter/        preface and opening pages
|-- chapters/           one file per chapter
|-- appendices/         one file per appendix
`-- figures/            images (PDF, PNG, JPG)
\end{lstlisting}

\begin{note}
Only \file{main.tex} and \file{style/booktemplate.sty} contain configuration.
Everything else is content that you are expected to edit or replace.
\end{note}

\section{Compiling the Book}\label{sec:compiling}

\subsection{On Overleaf}

\begin{enumerate}
  \item Download the template as a \file{.zip} archive.
  \item In Overleaf choose \emph{New Project} and then \emph{Upload Project},
        and select the archive.
  \item Open the \emph{Menu}, set the compiler to \emph{pdfLaTeX} and the main
        document to \file{main.tex}.
  \item Click \emph{Recompile}. Overleaf runs BibTeX automatically.
\end{enumerate}

\subsection{On Your Computer}

Install a \TeX{} distribution such as \TeX{} Live or MiKTeX and run:

\begin{lstlisting}[language={}]
latexmk main.tex
\end{lstlisting}

\noindent The file \file{.latexmkrc} tells \pkg{latexmk} to use pdfLaTeX and to
run BibTeX when needed.

\section{Adding a Chapter}\label{sec:adding-chapter}

Adding a chapter takes two steps. First, create a new file in the
\file{chapters/} folder:

\begin{lstlisting}[language={[LaTeX]TeX}]
% chapters/05-my-chapter.tex
\chapter{My Chapter}\label{chap:my-chapter}

Your text goes here.
\end{lstlisting}

\noindent Second, add one line to \file{main.tex}, in the position where the
chapter should appear:

\begin{lstlisting}[language={[LaTeX]TeX}]
\include{chapters/05-my-chapter}
\end{lstlisting}

\begin{tip}
While writing, uncomment the \cmd{includeonly} line in \file{main.tex} to
compile only the chapter you are working on. Page numbers, cross-references
and the contents are preserved from the previous full compilation.
\end{tip}
